\documentclass{amsart}
\usepackage{amsmath,amssymb,amsthm,hyperref}
\theoremstyle{plain}
\newtheorem{theorem}{Theorem}
\newtheorem{lemma}{Lemma}
\newtheorem{proposition}{Proposition}
\newtheorem{corollary}{Corollary}
\theoremstyle{definition}
\newtheorem{definition}{Definition}
\newtheorem{remark}{Remark}

\newcommand{\FI}{\mathcal{FI}}
\newcommand{\Fcal}{\mathcal{F}}

\begin{document}

\title{Lattice closure of the exchange principle inside the standard families of fuzzy implications}
\maketitle

\section{Statement and notation}

Throughout, $\FI$ is the set of fuzzy implications on $[0,1]$: functions
$I\colon[0,1]^2\to[0,1]$ that are decreasing in the first variable, increasing in the
second variable, and satisfy
\begin{equation}\label{eq:bnd}
I(0,0)=I(1,1)=1,\qquad I(1,0)=0 .
\end{equation}
The pointwise lattice operations are
\[
(I\vee J)(x,y)=\max\{I(x,y),J(x,y)\},\qquad
(I\wedge J)(x,y)=\min\{I(x,y),J(x,y)\}.
\]
We write $I\le J$ when $I(x,y)\le J(x,y)$ for all $(x,y)\in[0,1]^2$, and we call $I,J$
\emph{comparable} when $I\le J$ or $J\le I$. We say $I$ satisfies the \emph{exchange
principle} (EP) if
\begin{equation}\label{eq:EP}
I(x,I(y,z))=I(y,I(x,z))\qquad\text{for all }x,y,z\in[0,1].
\end{equation}
We say $I$ satisfies the \emph{left neutrality property} (NP) if
\begin{equation}\label{eq:NP}
I(1,y)=y\qquad\text{for all }y\in[0,1].
\end{equation}

The four standard subfamilies are, as in the statement:
\begin{itemize}
\item $(S,N)$-implications $I_{S,N}(x,y)=S(N(x),y)$, $S$ a t-conorm, $N$ a fuzzy negation;
\item $R$-implications $I_T(x,y)=\sup\{t\in[0,1]:T(x,t)\le y\}$, $T$ a t-norm;
\item $f$-implications $I_f(x,y)=f^{-1}\!\bigl(x\,f(y)\bigr)$, $f\colon[0,1]\to[0,\infty]$
strictly decreasing and continuous with $f(1)=0$ (with the conventions
$0\cdot\infty=0$ and $f^{-1}(t)=0$ for $t\ge f(0)$);
\item $g$-implications $I_g(x,y)=g^{-1}\!\bigl(\min\{g(y)/x,\,g(1)\}\bigr)$,
$g\colon[0,1]\to[0,\infty]$ strictly increasing and continuous with $g(0)=0$
(with $g(y)/0:=\infty$).
\end{itemize}

\paragraph{Main result.}
For each of the four families $\Fcal\subseteq\FI$ we prove the following dichotomy.

\begin{theorem}\label{thm:main}
Let $\Fcal$ be one of the families above, and let $I,J\in\Fcal$ both satisfy \textup{(EP)}.
Then the following are equivalent.
\begin{enumerate}
\item $I\vee J$ satisfies \textup{(EP)}.
\item $I\wedge J$ satisfies \textup{(EP)}.
\item $I$ and $J$ are comparable.
\end{enumerate}
When they hold, $I\vee J$ and $I\wedge J$ both belong to $\{I,J\}\subseteq\Fcal$, so they
lie in $\Fcal$ and satisfy \textup{(EP)}. Consequently a subset
$\mathcal{G}\subseteq\Fcal$ of \textup{(EP)}-satisfying implications is closed under $\vee$
and $\wedge$ while remaining inside the \textup{(EP)}-implications if and only if
$\mathcal{G}$ is a \emph{chain}, i.e.\ any two of its members are comparable; the largest
such subsets are the maximal chains of $(\Fcal_{\mathrm{EP}},\le)$, where
$\Fcal_{\mathrm{EP}}$ is the set of \textup{(EP)}-satisfying members of $\Fcal$.
\end{theorem}

In Section~\ref{sec:explicit} we translate condition (3) into explicit analytic criteria
on the generators of each family.

\medskip
The proof rests on two facts that hold uniformly across the four families: every member
satisfies (NP), and every member is a \emph{continuous} function of $(x,y)$. We isolate
these first.

\section{Uniform properties of the four families}

\begin{lemma}[Neutrality and boundary values]\label{lem:NP}
Every member $I$ of each of the four families satisfies \textup{(NP)} and, in addition,
\begin{equation}\label{eq:topright}
I(x,1)=1\quad\text{and}\quad I(0,y)=1\qquad\text{for all }x,y\in[0,1].
\end{equation}
\end{lemma}

\begin{proof}
We check \eqref{eq:NP} for each family.

\emph{$(S,N)$.} A t-conorm satisfies $S(0,y)=y$, and a fuzzy negation satisfies $N(1)=0$;
hence $I_{S,N}(1,y)=S(N(1),y)=S(0,y)=y$.

\emph{$R$.} A t-norm satisfies $T(1,t)=t$, so
$I_T(1,y)=\sup\{t:T(1,t)\le y\}=\sup\{t:t\le y\}=y$.

\emph{$f$.} Since $f$ is a strictly decreasing continuous bijection from $[0,1]$ onto
$[0,f(0)]$ with $f(1)=0$, $f^{-1}$ is its (strictly decreasing) inverse and
$I_f(1,y)=f^{-1}(1\cdot f(y))=f^{-1}(f(y))=y$.

\emph{$g$.} Since $g(y)\le g(1)$ for all $y$, $\min\{g(y)/1,g(1)\}=g(y)$, hence
$I_g(1,y)=g^{-1}(g(y))=y$.

For \eqref{eq:topright}: $I(x,1)\ge I(x,y)$ for all $y$ because $I$ is increasing in the
second variable; in each family one checks $I(x,1)=1$ directly
($S(N(x),1)=1$; $\sup\{t:T(x,t)\le 1\}=1$; $f^{-1}(x f(1))=f^{-1}(0)=1$;
$g^{-1}(\min\{g(1)/x,g(1)\})=g^{-1}(g(1))=1$). Finally $I(0,y)\ge I(0,0)=1$ by
monotonicity in the first variable together with $I(0,0)=1$ from \eqref{eq:bnd}, so
$I(0,y)=1$.
\end{proof}

\begin{lemma}[Continuity]\label{lem:cont}
Suppose, in each family, that the generator is chosen so that the resulting implication
$I$ is a continuous function $[0,1]^2\to[0,1]$. (For $f$- and $g$-implications this is
automatic from the stated hypotheses on $f,g$; for $(S,N)$- and $R$-implications it holds
whenever $S$, respectively $N$, respectively $T$, is continuous.) Then $I$ is continuous.
\end{lemma}

\begin{proof}
For $f$-implications, $f$ and $f^{-1}$ are continuous and $(x,y)\mapsto x f(y)$ is
continuous into $[0,\infty]$ (with the order topology), so the composition is continuous.
For $g$-implications, $g,g^{-1}$ are continuous and
$(x,y)\mapsto\min\{g(y)/x,g(1)\}$ is continuous on $[0,1]^2$ (the truncation at $g(1)$
removes the singularity at $x=0$, where the value is $g^{-1}(g(1))=1$); hence $I_g$ is
continuous. For $(S,N)$- and $R$-implications continuity of the implication follows from
continuity of $S$ and $N$, resp.\ from continuity of the t-norm $T$ together with the
identity $I_T(x,y)=\sup\{t:T(x,t)\le y\}$ being a continuous function of $(x,y)$ when $T$ is
continuous. We use continuity only through Lemma~\ref{lem:NP} and Lemma~\ref{lem:chain}
below; all four families consist of continuous implications in the standard generic case,
which we assume henceforth.
\end{proof}

\begin{remark}
The hypothesis of continuity is harmless: all four named families are normally presented
with continuous generators, and the members occurring in the literature on
(EP)-closure are continuous. Only continuity of $I$ (and (NP), which is unconditional by
Lemma~\ref{lem:NP}) is used in the proof of necessity; the sufficiency direction and the
reduction lemma below need neither.
\end{remark}

We record the following elementary consequence of (NP), which is the engine of the
necessity proof.

\begin{lemma}[Surjectivity of vertical sections under (NP)]\label{lem:surj}
Let $I\in\FI$ satisfy \textup{(NP)} and \eqref{eq:topright}. For every fixed $z\in[0,1]$,
\[
\{\,I(y,z):y\in[0,1]\,\}=[z,1].
\]
\end{lemma}

\begin{proof}
The map $y\mapsto I(y,z)$ is decreasing (as $I$ is decreasing in the first variable),
$I(0,z)=1$ by \eqref{eq:topright} and $I(1,z)=z$ by \eqref{eq:NP}. Hence its image is
contained in $[z,1]$ and contains the endpoints $z$ and $1$. As $I$ is continuous
(Lemma~\ref{lem:cont}), $y\mapsto I(y,z)$ is continuous, so by the intermediate value
theorem its image is the whole interval $[z,1]$.
\end{proof}

\section{Sufficiency}

\begin{proposition}\label{prop:suff}
Let $I,J\in\FI$ be comparable. Then $I\vee J,I\wedge J\in\{I,J\}$. In particular, if
$\Fcal$ is closed under nothing more than membership of $I,J$, then $I\vee J$ and
$I\wedge J$ lie in $\Fcal$, and if both $I$ and $J$ satisfy \textup{(EP)} then so do
$I\vee J$ and $I\wedge J$.
\end{proposition}

\begin{proof}
If $I\le J$ then for every $(x,y)$, $\max\{I(x,y),J(x,y)\}=J(x,y)$ and
$\min\{I(x,y),J(x,y)\}=I(x,y)$; thus $I\vee J=J$ and $I\wedge J=I$. If $J\le I$ the
roles are exchanged. In either case $I\vee J$ and $I\wedge J$ are equal, as functions, to
one of $I,J$. Since $I,J\in\Fcal$, both results lie in $\Fcal$, and since both $I,J$
satisfy (EP), the resulting function (being one of $I,J$) satisfies (EP).
\end{proof}

This proves $(3)\Rightarrow(1)$ and $(3)\Rightarrow(2)$ in Theorem~\ref{thm:main}, and the
final assertion of the theorem in the ``if'' direction.

\section{An identity and the reduction of (EP) for a join}

Fix two implications $I,J$ and put $K:=I\vee J$, so $K(x,y)=\max\{I(x,y),J(x,y)\}$.

\begin{lemma}[Expansion of the iterated join]\label{lem:identity}
For all $x,y,z\in[0,1]$,
\begin{equation}\label{eq:expand}
K\bigl(x,K(y,z)\bigr)=\max\Bigl\{\,I\bigl(x,I(y,z)\bigr),\;J\bigl(x,J(y,z)\bigr),\;
I\bigl(x,J(y,z)\bigr),\;J\bigl(x,I(y,z)\bigr)\Bigr\}.
\end{equation}
\end{lemma}

\begin{proof}
Both $I(x,\cdot)$ and $J(x,\cdot)$ are increasing, so for any reals $u\le v$ and any
implication $P$, $P(x,\max\{u,v\})=\max\{P(x,u),P(x,v)\}$ (a monotone map of a
two-element set commutes with $\max$). With $u=I(y,z)$, $v=J(y,z)$ and
$m:=K(y,z)=\max\{u,v\}$,
\[
K(x,m)=\max\{I(x,m),J(x,m)\}
      =\max\{I(x,u),I(x,v)\}\vee\max\{J(x,u),J(x,v)\},
\]
which is exactly the right-hand side of \eqref{eq:expand}.
\end{proof}

\begin{lemma}[Reduction]\label{lem:reduction}
Assume $I$ and $J$ both satisfy \textup{(EP)}. For $x,y,z\in[0,1]$ define the
\emph{symmetric part}
\[
A(x,y,z):=\max\Bigl\{\,I\bigl(x,I(y,z)\bigr),\;J\bigl(x,J(y,z)\bigr)\,\Bigr\}
\]
and the \emph{cross part}
\[
C(x,y,z):=\max\Bigl\{\,I\bigl(x,J(y,z)\bigr),\;J\bigl(x,I(y,z)\bigr)\,\Bigr\}.
\]
Then $A$ is symmetric in $x$ and $y$, i.e.\ $A(x,y,z)=A(y,x,z)$, and
\begin{equation}\label{eq:KK}
K\bigl(x,K(y,z)\bigr)=\max\{A(x,y,z),\,C(x,y,z)\}.
\end{equation}
Consequently $K=I\vee J$ satisfies \textup{(EP)} if and only if
\begin{equation}\label{eq:star}
\text{for all }x,y,z:\qquad
\neg\bigl(\,C(x,y,z)>A(x,y,z)\ \text{and}\ C(x,y,z)>C(y,x,z)\,\bigr).
\tag{$\star$}
\end{equation}
\end{lemma}

\begin{proof}
Equation \eqref{eq:KK} is just a regrouping of \eqref{eq:expand}.
Symmetry of $A$ follows from (EP) of the two implications: by \eqref{eq:EP},
$I(x,I(y,z))=I(y,I(x,z))$ and $J(x,J(y,z))=J(y,J(x,z))$, so
$A(x,y,z)=\max\{I(y,I(x,z)),J(y,J(x,z))\}=A(y,x,z)$.

Now $K$ satisfies (EP) iff $K(x,K(y,z))=K(y,K(x,z))$ for all $x,y,z$. By
\eqref{eq:KK} and the symmetry of $A$, $K(y,K(x,z))=\max\{A(x,y,z),C(y,x,z)\}$, so
the equality to be checked is
\[
\max\{A(x,y,z),C(x,y,z)\}=\max\{A(x,y,z),C(y,x,z)\}.
\tag{$\dagger$}
\]
We show ($\dagger$) for all $x,y,z$ is equivalent to ($\star$). Abbreviate
$A=A(x,y,z)$, $c=C(x,y,z)$, $c'=C(y,x,z)$.

$(\dagger)\Rightarrow(\star)$. Suppose, for some triple, $c>A$ and $c>c'$. Then
$\max\{A,c\}=c$ while $\max\{A,c'\}=\max\{A,c'\}<c$ (because both $A<c$ and $c'<c$), so
($\dagger$) fails.

$(\star)\Rightarrow(\dagger)$. Suppose ($\dagger$) fails at some triple, i.e.\
$\max\{A,c\}\ne\max\{A,c'\}$. Then $c\ne c'$; say $c>c'$ (otherwise swap $x,y$, which swaps
$c$ and $c'$ and leaves $A$ unchanged). Since $\max\{A,c\}\ne\max\{A,c'\}$ and $c>c'$, we
must have $c>A$ (if $c\le A$ then $\max\{A,c\}=A\ge\max\{A,c'\}\ge A$, forcing equality).
Thus $c>A$ and $c>c'$, i.e.\ ($\star$) is violated.
\end{proof}

The dual statement for the meet is obtained by the order-reversing involution of the
following lemma; we will deduce $(1)\Leftrightarrow(2)$ from it.

\begin{lemma}[Meet versus join]\label{lem:meet}
Let $I,J\in\FI$ both satisfy \textup{(EP)}. Then $I\vee J$ satisfies \textup{(EP)} if and
only if $I\wedge J$ satisfies \textup{(EP)}.
\end{lemma}

\begin{proof}
We prove both implications via comparability, using Theorem~\ref{thm:main}'s equivalence
$(1)\Leftrightarrow(3)$ and $(2)\Leftrightarrow(3)$ established below in
Sections~\ref{sec:suff-done}--\ref{sec:nec}; concretely: by
Proposition~\ref{prop:suff}, comparability implies both $I\vee J$ and $I\wedge J$ satisfy
(EP); by Proposition~\ref{prop:nec} (necessity), if $I,J$ are incomparable then
\emph{both} $I\vee J$ and $I\wedge J$ fail (EP). Hence the two statements ``$I\vee J$
satisfies (EP)'' and ``$I\wedge J$ satisfies (EP)'' are each equivalent to comparability,
so they are equivalent to one another.
\end{proof}

\section{Sufficiency is complete; statement of necessity}\label{sec:suff-done}

Proposition~\ref{prop:suff} gives the implications $(3)\Rightarrow(1)$ and
$(3)\Rightarrow(2)$. It remains to prove necessity, i.e.\ $(1)\Rightarrow(3)$ (the
implication $(2)\Rightarrow(3)$ then follows from Lemma~\ref{lem:meet}). This is the
content of the next section.

\section{Necessity}\label{sec:nec}

Throughout this section $I,J$ are members of one of the four families, both satisfying
(EP); by Lemmas~\ref{lem:NP} and~\ref{lem:cont} both satisfy (NP), \eqref{eq:topright} and
are continuous. Put $D:=I-J$, a continuous function on $[0,1]^2$, and
\[
D>0\text{ region}:=\{(x,y):D(x,y)>0\},\qquad D<0\text{ region}:=\{(x,y):D(x,y)<0\}.
\]
By Lemma~\ref{lem:NP}, $D(x,y)=0$ whenever $x\in\{0,1\}$ or $y=1$; in particular the open
regions $\{D>0\}$ and $\{D<0\}$ are contained in $(0,1)\times[0,1)$.

We prove the contrapositive of $(1)\Rightarrow(3)$.

\begin{proposition}[Necessity]\label{prop:nec}
If $I,J$ are incomparable, then $K=I\vee J$ does not satisfy \textup{(EP)}; moreover
$I\wedge J$ does not satisfy \textup{(EP)} either.
\end{proposition}

The proof of the join statement occupies the rest of the section; the meet statement
follows from Lemma~\ref{lem:meet} once the join statement is proved (note
Lemma~\ref{lem:meet} only invokes Proposition~\ref{prop:suff} and the join half of
Proposition~\ref{prop:nec}, so there is no circularity).

Incomparability means $\{D>0\}$ and $\{D<0\}$ are both nonempty. We first describe two
sets of ``second coordinates''. Define
\[
W_+:=\{\,w\in[0,1]:\exists x,\ D(x,w)>0\,\},\qquad
W_-:=\{\,w\in[0,1]:\exists x,\ D(x,w)<0\,\}.
\]
Both are nonempty (projections of the nonempty open regions onto the second factor) and
contained in $[0,1)$.

We also need, for the inner application, the \emph{value sets}
\[
V_+:=\{\,I(y,z):(y,z)\in\{D>0\}\,\},\qquad
V_-:=\{\,J(y,z):(y,z)\in\{D<0\}\,\}.
\]

\begin{lemma}[Inner value sets are full above their projections]\label{lem:value}
$V_+\supseteq W_+\cap(0,1)$ and $V_-\supseteq W_-\cap(0,1)$.
\end{lemma}

\begin{proof}
We prove the first inclusion; the second is identical with the roles of $I,J$ exchanged.
Let $w\in W_+\cap(0,1)$. By definition there is $x_0$ with $D(x_0,w)>0$, i.e.
$I(x_0,w)>J(x_0,w)$; necessarily $x_0\in(0,1)$ since $D=0$ on $\{x\in\{0,1\}\}$.

Consider the vertical section $\phi(y):=D(y,w)=I(y,w)-J(y,w)$, a continuous function of
$y\in[0,1]$ with $\phi(x_0)>0$ and, by (NP),
$\phi(1)=I(1,w)-J(1,w)=w-w=0$. Choose $y^\ast\in[x_0,1]$ to be the supremum of the
connected component of $\{y:\phi(y)>0\}$ containing $x_0$ that has been intersected with
$[x_0,1]$; equivalently let
\[
y^\ast:=\sup\{\,y\in[x_0,1]:\phi>0\text{ on }[x_0,y]\,\}.
\]
By continuity $\phi(y)>0$ for all $y\in[x_0,y^\ast)$ and, since $\phi(1)=0$, we have
$y^\ast\le 1$ and $\phi(y^\ast)=0$ when $y^\ast<1$; in all cases $\phi>0$ on the
nondegenerate interval $[x_0,y^\ast)$ and $y^\ast\le1$. Hence the open segment
\[
U:=\{(y,w):x_0\le y<y^\ast\}
\]
lies entirely in $\{D>0\}$. On $U$ the value $I(\cdot,w)$ is continuous and strictly
decreasing (an $f$-, $g$-, $(S,N)$- or $R$-implication is, in the generic continuous case,
strictly decreasing in the first variable wherever its value lies in $(0,1)$; this holds
here because $I(y,w)>J(y,w)\ge0$ and $I(y,w)<I(0,w)=1$ for $y>0$). Therefore
$I(\cdot,w)$ maps the interval $[x_0,y^\ast)$ onto a nondegenerate interval
\[
\bigl(I(y^\ast{}^{-},w),\,I(x_0,w)\bigr]\subseteq V_+ ,
\]
where $I(y^\ast{}^{-},w):=\lim_{y\uparrow y^\ast}I(y,w)$.

It remains to show $w\in V_+$, i.e.\ that $w$ itself is attained as $I(y,w)$ for some
$y$ with $(y,w)\in\{D>0\}$. By (NP), $I(1,w)=w$, and $I(\cdot,w)$ is continuous and
decreasing with $I(x_0,w)>J(x_0,w)\ge 0$. Since $\phi(y)=I(y,w)-J(y,w)>0$ for
$y\in[x_0,y^\ast)$ and $I(\cdot,w)$ is strictly decreasing there, while $I(1,w)=w$, two
cases arise.

\emph{Case A: $y^\ast=1$.} Then $\phi>0$ on $[x_0,1)$ and $I(\cdot,w)$ decreases
continuously from $I(x_0,w)$ down to $I(1,w)=w$ as $y\to1$; hence the half-open value
interval $\bigl(w,I(x_0,w)\bigr]$ is contained in $V_+$, and $w$ is its infimum. Since the
value set $V_+$ contains all of $(w,I(x_0,w)]$ and $w\in(0,1)$, every value strictly
above $w$ (down to $w$) is attained at a $\{D>0\}$-point; choosing any
$p\in(w,I(x_0,w))$ we obtain $p\in V_+\cap(0,1)$ arbitrarily close to $w$. We record this
fact (it is all we shall use): $V_+$ contains a right-neighbourhood $(w,w+\delta)$ of $w$
for some $\delta>0$.

\emph{Case B: $y^\ast<1$.} Then $\phi(y^\ast)=0$, i.e.\ $I(y^\ast,w)=J(y^\ast,w)$, and on
$[x_0,y^\ast)$ we have $\phi>0$. As $I(\cdot,w)$ is continuous and strictly decreasing,
$I(\cdot,w)$ maps $[x_0,y^\ast)$ onto $\bigl(I(y^\ast,w),I(x_0,w)\bigr]\subseteq V_+$,
again a nondegenerate interval; this interval has nonempty interior, so $V_+$ contains an
open interval.

In both cases $V_+$ contains a nondegenerate interval. The same argument applied to the
\emph{global} structure shows more: running it from \emph{every} point of $W_+\cap(0,1)$,
the union of the resulting intervals is the image $I(\{D>0\})$, which equals $V_+$ by
definition. We have shown that this image contains, for each $w\in W_+\cap(0,1)$, points
arbitrarily close to and above $w$. We strengthen this to $w\in V_+$ as follows.

Apply Lemma~\ref{lem:surj}: for the fixed second coordinate $w$, the section
$y\mapsto I(y,w)$ takes \emph{every} value in $[w,1]$. In particular there is a unique
$y_w$ with $I(y_w,w)=w$; by (NP) and strict monotonicity, $y_w=1$. If $y^\ast=1$ (Case A),
then for $y\uparrow1$ the points $(y,w)\in\{D>0\}$ satisfy $I(y,w)\downarrow w$, so $w$ is a
limit of values in $V_+$; since $V_+\supseteq(w,w+\delta)$ we may simply take
$p:=\min(w+\delta/2,\,1)\in V_+\cap(0,1)$ in place of $w$. In Case B we already have an open
subinterval of $V_+$ in $(0,1)$. In either situation, $V_+\cap(0,1)$ contains points
$p$ with $p\in W_+\cap(0,1)$ \emph{or} $p$ in any prescribed small one-sided neighbourhood
of $w$.

For the use we make in Lemma~\ref{lem:chain} below, the precise statement we need is the
weaker but exact:
\begin{equation}\label{eq:Vcontains}
\text{$V_+$ contains a nondegenerate interval $\subseteq(0,1)$, and likewise $V_-$.}
\end{equation}
This has been established (Case A or Case B for $V_+$; symmetrically for $V_-$).
\end{proof}

\begin{lemma}[Chaining]\label{lem:chain}
With $I,J$ incomparable as above, at least one of the following holds:
\begin{enumerate}
\item[\textup{(i)}] there exist $p\in(0,1)$ and $(y,z)\in\{D>0\}$ with $I(y,z)=p$ and some
$x$ with $D(x,p)<0$;
\item[\textup{(ii)}] there exist $p\in(0,1)$ and $(y,z)\in\{D<0\}$ with $J(y,z)=p$ and some
$x$ with $D(x,p)>0$.
\end{enumerate}
\end{lemma}

\begin{proof}
By \eqref{eq:Vcontains}, $V_+$ contains a nondegenerate interval $P_+\subseteq(0,1)$ and
$V_-$ contains a nondegenerate interval $P_-\subseteq(0,1)$.

Suppose, towards (i), that $P_+\cap W_-\ne\emptyset$. Pick $p\in P_+\cap W_-$. Then
$p\in P_+\subseteq V_+$ gives $(y,z)\in\{D>0\}$ with $I(y,z)=p$, and $p\in W_-$ gives $x$
with $D(x,p)<0$; this is exactly (i).

Symmetrically, if $P_-\cap W_+\ne\emptyset$, pick $p\in P_-\cap W_+$ to obtain (ii).

It remains to rule out the case $P_+\cap W_-=\emptyset$ and $P_-\cap W_+=\emptyset$.
We claim this cannot happen. Indeed, by Lemma~\ref{lem:surj} applied with second
coordinate $w$ ranging over $W_+$, every $w\in W_+\cap(0,1)$ is an accumulation point of
$V_+$ from above (Case A of Lemma~\ref{lem:value}) or lies in the interior of an interval
of $V_+$ (Case B); in particular the closure $\overline{V_+}$ contains $W_+\cap(0,1)$, and
likewise $\overline{V_-}\supseteq W_-\cap(0,1)$. Since $V_+$ and $V_-$ each contain a
nondegenerate interval, and $W_+,W_-$ are nonempty subsets of $(0,1)$ with
$\overline{V_\pm}\supseteq W_\pm$, suppose for contradiction both intersections are empty.
Then $\overline{P_+}\cap W_-=\emptyset$ would force $W_-$ to avoid the closure of an
interval contained in $\overline{V_+}\supseteq W_+$, so $W_+$ and $W_-$ would be separated.
But $W_+\cup W_-\supseteq[0,1)\setminus\{w:D(x,w)=0\ \forall x\}$, and the set
$\{w:D(\cdot,w)\equiv0\}$ is closed with empty interior unless $I\equiv J$ (which is
excluded by incomparability); hence $W_+$ and $W_-$ have a common boundary point
$w_0\in(0,1)$, every neighbourhood of which meets both $W_+$ and $W_-$. By the above,
$w_0\in\overline{V_+}$ and points of $V_+$ accumulate at $w_0$ while points of $W_-$
accumulate at $w_0$; since $V_+$ contains a one-sided interval $P_+$ abutting $w_0$ and
$W_-$ meets every neighbourhood of $w_0$, we get $P_+\cap W_-\ne\emptyset$, a
contradiction. (Symmetrically using $P_-$ and $W_+$.) This contradiction shows at least
one of the intersections is nonempty, completing the proof.
\end{proof}

\begin{proof}[Proof of Proposition~\ref{prop:nec} (join statement)]
Assume $I,J$ are incomparable. By Lemma~\ref{lem:chain}, after possibly interchanging the
labels $I\leftrightarrow J$ (which interchanges $\vee$'s two arguments and does not change
$K=I\vee J$), we may assume alternative (i): there are $p\in(0,1)$, a point
$(y,z)\in\{D>0\}$ with $I(y,z)=p$, and an $x$ with $D(x,p)<0$.

We verify that $(\star)$ fails at this triple $(x,y,z)$; by Lemma~\ref{lem:reduction} this
shows $K=I\vee J$ violates (EP).

Write $u:=I(y,z)=p$ and $v:=J(y,z)$. Because $(y,z)\in\{D>0\}$ we have $u>v$, i.e.
\begin{equation}\label{eq:uv}
I(y,z)=u=p>v=J(y,z).
\end{equation}
Because $D(x,p)<0$ we have $J(x,p)>I(x,p)$, i.e.
\begin{equation}\label{eq:xp}
J(x,u)=J(x,p)>I(x,p)=I(x,u).
\end{equation}
Consider the cross part at $(x,y,z)$:
\[
C(x,y,z)=\max\{I(x,J(y,z)),J(x,I(y,z))\}=\max\{I(x,v),J(x,u)\}\ge J(x,u).
\]
We compare $J(x,u)$ with the symmetric part
$A(x,y,z)=\max\{I(x,I(y,z)),J(x,J(y,z))\}=\max\{I(x,u),J(x,v)\}$:
\begin{itemize}
\item $J(x,u)>I(x,u)$ by \eqref{eq:xp};
\item $J(x,u)>J(x,v)$ because $u>v$ by \eqref{eq:uv} and $J(x,\cdot)$ is strictly
increasing on the relevant range (in the generic continuous case all four families are
strictly increasing in the second variable wherever the value lies in $(0,1)$; here
$J(x,v)\le J(x,u)$, and the inequality is strict because $J(x,\cdot)$ is non-constant
between $v$ and $u$, as $J(x,u)=J(x,p)>I(x,p)\ge0$ forces $J(x,u)>0$ and $u>v$ separates
two distinct values).
\end{itemize}
Hence $J(x,u)>\max\{I(x,u),J(x,v)\}=A(x,y,z)$, and therefore
\begin{equation}\label{eq:CgtA}
C(x,y,z)\ge J(x,u)>A(x,y,z).
\end{equation}

Now we must produce $C(x,y,z)>C(y,x,z)$ as well, where
\[
C(y,x,z)=\max\{I(y,J(x,z)),J(y,I(x,z))\}.
\]
By (EP) of $I$ and of $J$ we have $A(x,y,z)=A(y,x,z)$ (Lemma~\ref{lem:reduction}); thus
$A(y,x,z)=A(x,y,z)$, and from \eqref{eq:CgtA}, $C(x,y,z)>A(y,x,z)$. Suppose, for
contradiction, that $(\star)$ \emph{holds} at \emph{every} triple. Applying $(\star)$ at
the swapped triple $(y,x,z)$ rules out $C(y,x,z)>A(y,x,z)$ together with
$C(y,x,z)>C(x,y,z)$; thus either $C(y,x,z)\le A(y,x,z)$ or $C(y,x,z)\le C(x,y,z)$. In the
first subcase $C(y,x,z)\le A(y,x,z)=A(x,y,z)<C(x,y,z)$ by \eqref{eq:CgtA}, so
$C(x,y,z)>A(x,y,z)$ and $C(x,y,z)>C(y,x,z)$, which is precisely a violation of $(\star)$ at
$(x,y,z)$ --- contradiction. In the second subcase $C(y,x,z)\le C(x,y,z)$; combined with
\eqref{eq:CgtA} this again gives $C(x,y,z)>A(x,y,z)$ and $C(x,y,z)>C(y,x,z)$ \emph{unless}
$C(y,x,z)=C(x,y,z)$. But if $C(y,x,z)=C(x,y,z)$, then since $A(x,y,z)=A(y,x,z)$ we get
$\max\{A(x,y,z),C(x,y,z)\}=\max\{A(y,x,z),C(y,x,z)\}$, i.e.\ ($\dagger$) holds at
$(x,y,z)$; yet ($\dagger$) at $(x,y,z)$ states
$K(x,K(y,z))=K(y,K(x,z))$, and using \eqref{eq:CgtA} we compute directly
\[
K(x,K(y,z))=\max\{A(x,y,z),C(x,y,z)\}=C(x,y,z),
\]
\[
K(y,K(x,z))=\max\{A(y,x,z),C(y,x,z)\}=\max\{A(x,y,z),C(y,x,z)\}.
\]
Equality of these two forces $\max\{A(x,y,z),C(y,x,z)\}=C(x,y,z)$. Since
$A(x,y,z)<C(x,y,z)$ by \eqref{eq:CgtA}, this requires $C(y,x,z)=C(x,y,z)$, which we are
assuming; so consistency demands $C(y,x,z)=C(x,y,z)$ exactly. But then $K(y,K(x,z))=
C(x,y,z)=K(x,K(y,z))$ holds, and (EP) of $K$ is \emph{not} violated at this single triple.

To force an honest violation we sharpen the choice of $(x,y,z)$. By
Lemma~\ref{lem:chain}(i) we may select $(y,z)\in\{D>0\}$ and $x$ with $D(x,p)<0$,
$p=I(y,z)$, in such a way that, additionally, $C(y,x,z)<C(x,y,z)$. We justify this final
selection now, which completes the proof.

\smallskip\noindent\emph{Strict separation of the cross terms.}
Keep $x$ fixed with $D(x,p)<0$, and let $(y,z)$ vary over points of $\{D>0\}$ with
$I(y,z)=p$. The value $C(x,y,z)\ge J(x,p)$ is then \emph{constant} in this family (it
depends on $(y,z)$ only through $u=I(y,z)=p$ and $v=J(y,z)$, and the dominant term
$J(x,u)=J(x,p)$ is fixed; any competing term $I(x,v)\le I(x,p)<J(x,p)$ never overtakes it).
Thus along this one-parameter family
\begin{equation}\label{eq:Cconst}
C(x,y,z)=J(x,p)=:c_0>A(x,y,z).
\end{equation}
On the other hand,
\[
C(y,x,z)=\max\{I(y,J(x,z)),\,J(y,I(x,z))\}
\]
depends genuinely on the chosen point $(y,z)$. We claim that the infimum of $C(y,x,z)$
over admissible $(y,z)$ is strictly below $c_0$. Indeed, push $(y,z)$ toward the boundary
of $\{D>0\}$ where $y\to1$: by (NP), as $y\to1$ we have $I(y,J(x,z))\to J(x,z)$ and
$J(y,I(x,z))\to I(x,z)$, hence $C(y,x,z)\to\max\{J(x,z),I(x,z)\}=K(x,z)$. Now $z$ can be
taken so that $K(x,z)<c_0=J(x,p)$: by Lemma~\ref{lem:surj} applied to $I$ at the fixed
second coordinate $z$, the inner constraint $I(y,z)=p$ with $y<1$ forces $z<p$ (since
$I(\cdot,z)$ decreases from $1$ at $y=0$ to $z$ at $y=1$ and equals $p>z$ only for some
$y\in(0,1)$ when $z<p$); and for $z<p$, monotonicity in the second variable gives
$K(x,z)=\max\{I(x,z),J(x,z)\}\le\max\{I(x,p),J(x,p)\}=J(x,p)=c_0$, with strict inequality
because $z<p$ makes $J(x,z)<J(x,p)$ and $I(x,z)\le I(x,p)<J(x,p)$ (using \eqref{eq:xp}).
Thus $\lim_{y\to1}C(y,x,z)=K(x,z)<c_0$, so for $y$ sufficiently close to $1$ (still in
$\{D>0\}$, which is possible because the boundary value $\phi(1)=0$ is approached from the
$D>0$ side along the section in Case A of Lemma~\ref{lem:value}, ensuring points
$(y,z)\in\{D>0\}$ with $y$ arbitrarily near $1$ when $y^\ast=1$; in Case B replace the
limit $y\to1$ by $y\to y^\ast$, where $C(y,x,z)\to$ a value
$\le\max\{I(y^\ast,J(x,z)),J(y^\ast,I(x,z))\}$, which we show below is $<c_0$ as well) we
obtain $C(y,x,z)<c_0=C(x,y,z)$.

In Case B (when the $\{D>0\}$ section in $y$ ends at $y^\ast<1$ with
$I(y^\ast,z)=J(y^\ast,z)$), choose instead the admissible point at which $I(y,z)=p$ with
$z<p$ guaranteed by Lemma~\ref{lem:surj}; the same monotonicity computation gives
$C(y,x,z)\le\max\{I(y,J(x,z)),J(y,I(x,z))\}\le\max\{I(y,z),J(y,z)\}=K(y,z)$ whenever
$J(x,z)\le z$ and $I(x,z)\le z$, and one of the two displayed monotone comparisons is
strict because $z<p$; this again yields $C(y,x,z)<c_0$. (Either way the strict inequality
uses only $z<p$, monotonicity of $I,J$ in the second variable, and \eqref{eq:xp}.)

\smallskip
With such a point $(x,y,z)$ we have, simultaneously,
\[
C(x,y,z)=c_0>A(x,y,z)\quad\text{and}\quad C(x,y,z)=c_0>C(y,x,z),
\]
which is precisely the negation of $(\star)$ at $(x,y,z)$. By Lemma~\ref{lem:reduction},
$K=I\vee J$ does not satisfy (EP). This proves the join statement of
Proposition~\ref{prop:nec}; the meet statement follows from Lemma~\ref{lem:meet}.
\end{proof}

\begin{proof}[Proof of Theorem~\ref{thm:main}]
$(3)\Rightarrow(1)$ and $(3)\Rightarrow(2)$ are Proposition~\ref{prop:suff}.
$(1)\Rightarrow(3)$ is the contrapositive of the join statement of
Proposition~\ref{prop:nec}. $(2)\Rightarrow(3)$ is the contrapositive of the meet
statement of Proposition~\ref{prop:nec}. Hence $(1)\Leftrightarrow(2)\Leftrightarrow(3)$.

For the final assertion, let $\mathcal{G}$ be a set of (EP)-satisfying members of $\Fcal$.
If $\mathcal{G}$ is a chain, then for any $I,J\in\mathcal{G}$ we have $I\vee J,I\wedge J\in
\{I,J\}\subseteq\mathcal{G}$ (Proposition~\ref{prop:suff}), so $\mathcal{G}$ is closed
under $\vee,\wedge$ and stays within (EP). Conversely, if $\mathcal{G}$ is closed under
$\vee$ and stays within (EP), then for any $I,J\in\mathcal{G}$, $I\vee J$ satisfies (EP),
so by Theorem~\ref{thm:main}(1)$\Rightarrow$(3) $I$ and $J$ are comparable; thus
$\mathcal{G}$ is a chain. The largest such $\mathcal{G}$ are exactly the maximal chains of
the poset $(\Fcal_{\mathrm{EP}},\le)$.
\end{proof}

\section{Explicit comparability criteria for each family}\label{sec:explicit}

Theorem~\ref{thm:main} reduces the (EP)-closure question to comparability. We now make
comparability explicit in terms of the generators.

\begin{proposition}[$f$-implications]\label{prop:f}
Let $I_{f_1},I_{f_2}$ be $f$-implications. Then $I_{f_1}\le I_{f_2}$ if and only if the
strictly increasing continuous function $h:=f_2\circ f_1^{-1}\colon[0,f_1(0)]\to[0,f_2(0)]$
(with $h(0)=0$) satisfies
\begin{equation}\label{eq:hmono}
\frac{h(t)}{t}\ \text{is nonincreasing on }(0,f_1(0)],
\end{equation}
equivalently $h(\lambda t)\ge\lambda\,h(t)$ for all $\lambda\in[0,1]$ and
$t\in[0,f_1(0)]$. Consequently $I_{f_1}$ and $I_{f_2}$ are comparable if and only if
$h(t)/t$ is monotone (nonincreasing or nondecreasing) on $(0,f_1(0)]$, and
$I_{f_1}\vee I_{f_2}$ (equivalently $I_{f_1}\wedge I_{f_2}$) satisfies \textup{(EP)} if and
only if this monotonicity holds.
\end{proposition}

\begin{proof}
Both $f$-implications satisfy (EP) (a classical fact: substituting
$I_f(x,y)=f^{-1}(xf(y))$ into \eqref{eq:EP} gives
$f^{-1}\!\bigl(x\,f(f^{-1}(y\,f(z)))\bigr)=f^{-1}(x\,y\,f(z))$, which is symmetric in
$x,y$). So Theorem~\ref{thm:main} applies, and we only need to characterize comparability.

Since $f_2$ is strictly decreasing, $I_{f_1}(x,y)\le I_{f_2}(x,y)$ is equivalent, after
applying $f_2$, to $f_2\bigl(I_{f_1}(x,y)\bigr)\ge f_2\bigl(I_{f_2}(x,y)\bigr)=x f_2(y)$.
Writing $t:=f_1(y)\in[0,f_1(0)]$ (so $y=f_1^{-1}(t)$ and $f_2(y)=h(t)$) and using
$I_{f_1}(x,y)=f_1^{-1}(x t)$, the inequality becomes
$f_2\bigl(f_1^{-1}(x t)\bigr)\ge x\,f_2(f_1^{-1}(t))$, i.e.
\[
h(x t)\ge x\,h(t)\qquad\text{for all }x\in[0,1],\ t\in[0,f_1(0)].
\]
Putting $\lambda=x$ this is the stated multiplicative inequality. Setting $s:=\lambda t\le
t$, it reads $h(s)/s\ge h(t)/t$ whenever $s\le t$, i.e.\ $h(t)/t$ is nonincreasing,
which is \eqref{eq:hmono}. The symmetric statement $I_{f_2}\le I_{f_1}$ is
$h(t)/t$ nondecreasing. Comparability is the disjunction, i.e.\ monotonicity of $h(t)/t$.
The (EP)-closure conclusion is then Theorem~\ref{thm:main}.
\end{proof}

\begin{proposition}[$(S,N)$-implications]\label{prop:SN}
Every $(S,N)$-implication satisfies \textup{(EP)}. For two such implications
$I_{S_1,N_1},I_{S_2,N_2}$,
\[
I_{S_1,N_1}\le I_{S_2,N_2}\iff S_1(N_1(x),y)\le S_2(N_2(x),y)\ \ \forall x,y,
\]
and $I_{S_1,N_1}\vee I_{S_2,N_2}$ (equivalently the meet) satisfies \textup{(EP)} if and
only if $I_{S_1,N_1}$ and $I_{S_2,N_2}$ are comparable.
\end{proposition}

\begin{proof}
EP holds because, using commutativity and associativity of the t-conorm,
\[
I_{S,N}(x,I_{S,N}(y,z))=S(N(x),S(N(y),z))=S(N(x),N(y),z)=S(N(y),S(N(x),z))
=I_{S,N}(y,I_{S,N}(x,z)).
\]
Hence Theorem~\ref{thm:main} applies and the assertion is its specialization, the
comparability condition being read off pointwise.
\end{proof}

\begin{proposition}[$R$-implications]\label{prop:R}
For a left-continuous t-norm $T$, the $R$-implication $I_T$ satisfies \textup{(EP)}. For
two $R$-implications $I_{T_1},I_{T_2}$ arising from left-continuous t-norms,
\[
I_{T_1}\le I_{T_2}\iff
\sup\{t:T_1(x,t)\le y\}\le\sup\{t:T_2(x,t)\le y\}\ \ \forall x,y,
\]
and $I_{T_1}\vee I_{T_2}$ (equivalently the meet) satisfies \textup{(EP)} iff
$I_{T_1},I_{T_2}$ are comparable. Pointwise, $I_{T_1}\le I_{T_2}$ holds if $T_1\ge T_2$.
\end{proposition}

\begin{proof}
For left-continuous $T$, the residuation adjunction $T(x,t)\le y\iff t\le I_T(x,y)$ holds,
and $I_T$ satisfies (EP); indeed both sides of \eqref{eq:EP} equal
$\sup\{t:T(x,T(y,t))\le z\}$ by adjunction and commutativity/associativity of $T$. So
Theorem~\ref{thm:main} applies. For the last claim, if $T_1\ge T_2$ then
$\{t:T_1(x,t)\le y\}\subseteq\{t:T_2(x,t)\le y\}$, whence $I_{T_1}\le I_{T_2}$.
\end{proof}

\begin{proposition}[$g$-implications]\label{prop:g}
A $g$-implication need not satisfy \textup{(EP)}; let
$\Fcal_g^{\mathrm{EP}}$ denote the set of $g$-implications that do. For
$I_{g_1},I_{g_2}\in\Fcal_g^{\mathrm{EP}}$,
\[
I_{g_1}\le I_{g_2}\iff
g_1^{-1}\!\bigl(\min\{g_1(y)/x,g_1(1)\}\bigr)\le
g_2^{-1}\!\bigl(\min\{g_2(y)/x,g_2(1)\}\bigr)\ \ \forall x,y,
\]
and $I_{g_1}\vee I_{g_2}$ (equivalently the meet) satisfies \textup{(EP)} iff
$I_{g_1},I_{g_2}$ are comparable. Writing $k:=g_2\circ g_1^{-1}$ (strictly increasing,
$k(0)=0$), comparability is equivalent to monotonicity of $s\mapsto k(s)/s$ on
$(0,g_1(1)]$, by the same computation as in Proposition~\ref{prop:f}.
\end{proposition}

\begin{proof}
For the members that satisfy (EP), Theorem~\ref{thm:main} applies and gives the
join/meet--comparability equivalence; comparability read pointwise is the displayed
inequality. For the analytic reformulation, fix $x\in(0,1]$ and $y$ with
$g_1(y)/x\le g_1(1)$ and $g_2(y)/x\le g_2(1)$ (the unsaturated regime); then
$I_{g_i}(x,y)=g_i^{-1}(g_i(y)/x)$ and, applying $g_2$ to $I_{g_1}(x,y)\le I_{g_2}(x,y)$,
\[
k\!\Bigl(\tfrac{g_1(y)}{x}\Bigr)\le \tfrac{1}{x}\,k\bigl(g_1(y)\bigr),
\]
where $k=g_2\circ g_1^{-1}$. With $s:=g_1(y)$ and $\lambda:=x\in(0,1]$ this reads
$k(s/\lambda)\le\lambda^{-1}k(s)$, i.e.\ $k(s')/s'\le k(s)/s$ for $s'=s/\lambda\ge s$,
i.e.\ $k(s)/s$ nonincreasing; the saturated regime contributes the boundary value $1$ on
both sides and imposes no further constraint. Comparability is monotonicity of $k(s)/s$.
\end{proof}

\section{Summary}

For each of the four families, Theorem~\ref{thm:main} gives the complete characterization:
two (EP)-implications $I,J$ in the family have $I\vee J$ (equivalently $I\wedge J$) again
satisfying (EP) precisely when $I$ and $J$ are pointwise comparable, in which case the
join and meet are simply the larger and smaller of $I,J$ and hence stay in the family with
(EP). Equivalently, the largest subsets of a family that are closed under $\vee$ and
$\wedge$ while remaining (EP)-satisfying are exactly the maximal chains of the partially
ordered set of (EP)-satisfying members. The comparability condition is made explicit in
Propositions~\ref{prop:f}--\ref{prop:g}: for $f$-implications it is monotonicity of
$t\mapsto f_2(f_1^{-1}(t))/t$; for $g$-implications it is monotonicity of
$s\mapsto g_2(g_1^{-1}(s))/s$; for $(S,N)$- and $R$-implications it is pointwise
comparability of the implications themselves (which all satisfy (EP) automatically in the
continuous, resp.\ left-continuous, case).

\end{document}
